% Generated draft: requires mathematical and source review.
\documentclass[11pt]{article}
\usepackage[margin=25mm]{geometry}
\usepackage{fontspec}
\setmainfont{Latin Modern Roman}
\usepackage{amsmath,amssymb,mathtools,unicode-math}
\setmathfont{Latin Modern Math}
\usepackage{enumitem,needspace,xurl,hyperref}
\setlength{\parindent}{0pt}
\setlength{\parskip}{5pt}
\setlength{\emergencystretch}{3em}
\newcommand{\N}{\mathbb{N}}
\newcommand{\Z}{\mathbb{Z}}
\newcommand{\Q}{\mathbb{Q}}
\newcommand{\R}{\mathbb{R}}
\newcommand{\C}{\mathbb{C}}
\newcommand{\F}{\mathbb{F}}
\newcommand{\A}{\mathbb{A}}
\newcommand{\PP}{\mathbb{P}}
\newcommand{\E}{\mathbb{E}}
\newcommand{\eps}{\varepsilon}
\newcommand{\dd}{\,\mathrm{d}}
\begin{document}
\section*{UnsolvedMath ID: 11000048}
\textbf{Source problem code:} catalogue entry
\subsection{Definitions and mathematical statement}
Let \(\mathcal A_g\) be the moduli space of principally polarized complex abelian varieties of dimension \(g\), modeled by \(\operatorname{Sp}(2g,\Z)\backslash\mathfrak H_g\), and let \(\mathcal J_g\) denote the intended Schottky locus of Jacobians. Whether its closure, including degenerate/decomposable limits, is included is not specified here.

Using the preceding source's construction \(\operatorname{Cone}(\mathcal A_g)\), identified there with a \(g\)-dimensional Euclidean sector, describe explicitly the subset determined by \(\mathcal J_g\). In an asymptotic-cone interpretation, this means the cone points represented by appropriately scaled sequences lying in \(\mathcal J_g\), expressed in the stated sector coordinates. The metric, rescaling/basepoint convention, equivalence of sequences and sector coordinates are absent, so no particular cone construction or coordinate inequalities are invented. The source explains that the retained coordinates record changing relative sizes of basis vectors of tori while discarding skewing parameters, and hopes that removing those parameters simplifies the classical Schottky problem.

The continuation also proposes a coarse distortion question. For a subset \(S\) of a metric space \((Y,d_Y)\), compare the restricted distance \(d_{\mathrm{ext}}(x,y)=d_Y(x,y)\) with its induced path distance
\[
 d_{\mathrm{int}}(x,y)=\inf\{\operatorname{length}_{d_Y}(\gamma):
       \gamma\text{ is a path in }S\text{ joining }x\text{ to }y\},
\]
allowing \(+\infty\) when no rectifiable path exists. Apply this to the Schottky locus in \(\mathcal A_g\), with the intended metric and comparison convention. The source describes distortion as expected rather than established and also recalls Hain and Toledo's question of computing the second fundamental form of the Schottky locus. That differential-geometric question needs a chosen metric and smooth locus; neither computation nor a distortion growth rate is supplied.
\subsection{Short English statement}
Describe the coarse Schottky locus in the stated Euclidean-sector cone of abelian-variety moduli, and investigate its distortion by comparing ambient and intrinsic path distances.
\subsection{Sources}
\begin{itemize}
\item \emph{Problems on Mapping Class Groups and Related Topics}, edited by Farb, Problem 4.11, ``Coarse Schottky problem,'' and following distortion discussion, supplied PDF page 42. \url{https://www.math.uchicago.edu/~farb/papers/mcgbook.pdf}. Recorded access 2026-07-29; extraction pdf-font-structure. The cone definition and Subsection 3.2 distortion background are not supplied.
\item The statement attributes the second-fundamental-form problem to Hain and Toledo, perhaps among others, without a specific bibliography entry; that qualification is retained.
\item Local UnsolvedMath registry, \texttt{data/unsolvedmath/problems.json}, permanent ID \(11000048\), exact source code \texttt{catalogue entry}; supplied OPEN-TRIAGE summary reports unverified status.
\end{itemize}
\end{document}
